\begin{table*}[t]
\centering
\caption{Mean distance and F-score at 1 m under four post-processing and registration protocols. Means over three splits.}
\label{tab:protocol}
\footnotesize
\setlength{\tabcolsep}{3.6pt}
\begin{tabular}{@{}lrrrrrrrr@{}}
\toprule
 & \multicolumn{2}{c}{Mesh as extracted} & \multicolumn{2}{c}{+ tie-point alignment} & \multicolumn{2}{c}{+ Taubin smoothing} & \multicolumn{2}{c}{ICP registration} \\
\cmidrule(lr){2-3}\cmidrule(lr){4-5}\cmidrule(lr){6-7}\cmidrule(lr){8-9}
Method & Mean (m) & F@1 m & Mean (m) & F@1 m & Mean (m) & F@1 m & Mean (m) & F@1 m \\
\midrule
\multicolumn{9}{@{}l}{\textit{Itokawa (AMICA, phase 10\textdegree{})}} \\
2DGS & 6.58 & 0.532 & 6.55 & 0.540 & 6.55 & 0.546 & 6.56 & 0.539 \\
GOF & 2.41 & 0.365 & 2.40 & 0.365 & 2.40 & 0.362 & 2.39 & 0.373 \\
PGSR & \underline{1.18} & \underline{0.695} & 0.99 & \underline{0.811} & 0.99 & \underline{0.812} & 0.98 & \underline{0.821} \\
Asteroid-GS & 1.52 & 0.432 & 1.01 & 0.648 & 1.01 & 0.649 & 0.99 & 0.661 \\
Asteroid-GS$^{\dagger}$ & 1.53 & 0.415 & \underline{0.94} & 0.659 & \underline{0.94} & 0.658 & \underline{0.91} & 0.684 \\
\textbf{PCGS (this work)} & \textbf{0.59} & \textbf{0.853} & \textbf{0.59} & \textbf{0.854} & \textbf{0.59} & \textbf{0.853} & \textbf{0.59} & \textbf{0.857} \\
\addlinespace[3pt]
\multicolumn{9}{@{}l}{\textit{Ryugu (ONC-T, phase 19\textdegree{})}} \\
2DGS & 2.51 & \underline{0.388} & 2.51 & 0.410 & 2.50 & 0.416 & 2.48 & 0.421 \\
GOF & 4.20 & 0.156 & 4.21 & 0.177 & 4.20 & 0.162 & 3.38 & 0.255 \\
PGSR & \underline{1.51} & 0.383 & \underline{0.73} & \underline{0.788} & \underline{0.73} & \underline{0.787} & \underline{0.72} & \underline{0.794} \\
Asteroid-GS & 2.41 & 0.239 & 1.27 & 0.524 & 1.27 & 0.525 & 1.26 & 0.530 \\
Asteroid-GS$^{\dagger}$ & 2.43 & 0.242 & 1.29 & 0.529 & 1.29 & 0.530 & 1.27 & 0.539 \\
\textbf{PCGS (this work)} & \textbf{0.62} & \textbf{0.797} & \textbf{0.62} & \textbf{0.798} & \textbf{0.62} & \textbf{0.798} & \textbf{0.61} & \textbf{0.803} \\
\bottomrule
\end{tabular}
\tablefoot{Mesh as extracted: each method's own extraction, registered with the shared scene transform (for PCGS this already contains its offset to the SfM tie points). Tie-point alignment: translation and normal offset fitted to the tie points, no reference used. ICP registration: the mesh after tie-point alignment with a similarity transform (scale, rotation, translation) fitted to the reference by ICP instead of the shared scene transform, which removes any global misregistration of that mesh. $^{\dagger}$Asteroid-GS with the scene extent that its released code derives from the camera positions; the other rows use the object extent. Bold: best; underlined: second best.}
\end{table*}
